\chapter{Quietude e Resposta}

\lettrine{Q}{uando começámos o retiro} de Inverno, eu pedi-vos para aceitarem a totalidade
do que acontecer nos próximos dois meses. Estabeleçam a vossa intenção não
apenas para ter o tipo de retiro que gostariam, mas para se abrirem à
possibilidade de tudo o que surgir. Psicologicamente, isto prepara-nos para o
modo como a vida se move e muda. Quando programamos a nossa mente de forma a
tentar transformar a nossa vida naquilo que queremos, logo sempre nos sentimos
frustrados quando ela não corre da maneira que gostaríamos. Portanto, tentemos
mudar a nossa atitude para uma atitude de aceitação, vontade de olhar e
compreender as experiências, em vez de simplesmente tentarmos ver-nos livres
delas.

Estamos a tentar desenvolver esta prática de quietude. A quietude que está em
toda parte, estejamos num grupo ou sozinhos. Para estarmos com o silêncio, temos
que realizar a quietude, o silêncio. Em outras palavras, sermos dessa forma --
estarmos serenos e em silêncio.

Se simplesmente seguirmos as sensações agitadas do corpo e as proliferações da
mente, então claro está, o silêncio é impossível. Pode até ser uma experiência
ameaçadora, por a pessoa estar muito identificada com a agitação e inquietação
do reino sensorial, e incessantemente à procura de nascer nessa dimensão.

A ênfase agora é reconhecer essa inquietação pelo que ela é, não mais segui-la,
mas treinarmo-nos em direcção à serenidade. Não significa apenas suprimir a
forma física e oprimi-la, mas treiná-la, porque estes corpos precisam de ser
treinados com bondade. Se brutalizarmos os animais, eles não ficam muito
simpáticos, ou ficam? Eles tornam-se simplesmente criaturas assustadas,
desconfiadas e miseráveis. Treinar um animal não significa apenas mimá-lo, mas
orientá-lo. É o mesmo com o nosso próprio corpo. O nosso corpo precisa de ser
respeitado e orientado para não seguir a sua irrequieta energia e hábitos
inquietos.

Mas também não significa que devamos negar tudo. Um treinador precisa de ser
alguém amável e firme, não teimoso ou bruto. Não é ser amável no sentido de
ceder a tudo -- porque isso não é realmente ser amável -- mas cuidar,
preocupar-se, ter a quantidade certa de interesse. A atitude correcta em relação
ao seu próprio corpo e mente.

Como acalmar o corpo? Uma das formas é através da ``meditação de varrimento'', na
qual ``varremos'' o corpo com a nossa atenção, concentrando-nos nas sensações do
corpo enquanto o perfazemos. O corpo precisa de ser apercebido e aceite pelo que
é. Assim, trazemos à consciência até mesmo as tensões, sensações desagradáveis e
partes insensíveis do corpo. Ao fazermos isso, indo do topo da cabeça às plantas
dos pés e subindo novamente, o corpo sentir-se-á descontraído. É uma meditação
muito saudável e ajudará a treinar a mente para não ser apanhada na proliferação
conceitual e em divagações intermináveis.

Então, à medida que essas formações se começam a acalmar, começamos a sentir-nos
muito mais conscientes do silêncio da mente. Podemos permanecer nesse vazio cada
vez mais, onde não há eu, mas apenas o momento presente como é. A quietude e o
silêncio estão sempre presentes onde quer que estejamos, não importando em que
condição estamos.

Podemos permanecer no vazio ao simplesmente ficarmos de pé entre as árvores
desfolhadas do Inverno e olharmos para elas sem criarmos o que seja nelas.
Podemos sentir uma sensação de contentamento perfeito e calma perfeita, apenas
por se estar quieto e silencioso como as árvores. Talvez os nossos egos possam
dizer: ``bem, eu não quero tornar-me como uma árvore. Quero expressar a minha
verdadeira criatividade interior, a minha personalidade singular.'' Nós ouvimos
as vozes internas que reclamam e resmungam, o desejo de nos tornarmos algo,
aquilo que se evidencia ou existe. Mas não estamos a alimentar estas criaturas,
estamos a deixá-las ir e a mover-nos em direcção à quietude e ao silêncio.

Esta palavra ``existência'' significa ``apresentar-se''. Algo que não existe não
se apresenta. Então, quando dizemos ``não existente'', não quer dizer
matarmo-nos e não estarmos mais vivos, mas sim não seguir mais o desejo de nos
evidenciarmos, de nos tornarmos algo, de nos destacarmos. Ora, isso soa
realmente como uma perspectiva niilista: ``Ajahn Sumedho não quer existir! Oh,
pobre homem, precisa de ir a um psiquiatra.'' Mas não-existência não significa
que não queremos ter personalidade alguma, que nos queremos simplesmente tornar
pessoas taciturnas e desoladas. Não é isso. É a capacidade de permanecermos na
subtileza de estarmos simplesmente cientes, abertos e sensíveis, sem sermos
apanhados na ilusão de nos tornarmos outra coisa ou de destacarmo-nos de alguma
forma. É simplesmente realizar a paz da não-existência -- porque a
não-existência é pacífica. E quando há não-existência e vazio, há o
conhecimento, o brilho, a sabedoria, a consciência, a clareza, a iluminação. As
coisas são como são, aquilo que é, como é.

Nos valores ocidentais, a ênfase está em ser-se especial, um indivíduo único,
uma criança de Deus. Essa atitude é muito apoiada pela cultura e religião.
Existe o ``povo escolhido de Deus'', as seitas que sentem que foram chamadas por
Jesus (e todas as outras não); e são eles quem vão ter sucesso e viver num
paraíso eterno.

Mas com todas estas noções de se ser especial, de ser um indivíduo e todas essas
concepções egotistas -- o que é que acontece connosco? Por minha própria
experiência, o resultado de tudo isto foi sofrimento. Parecia existir um
tremendo investimento em ter uma personalidade única e vibrante. Às vezes eu
costumava pensar: ``Espera aí, talvez eu não tenha uma personalidade muito
simpática. Talvez eu não tenha personalidade.'' Havia tanta ansiedade,
frustração, ciúme e medo. Eu não queria ser um fracasso, não queria ser uma
mediocridade, ser uma ``pessoa qualquer''. Era muito doloroso estar sempre preso
naquele desejo de me tornar alguém. E enquanto uma pessoa tiver este desejo, vai
sempre ter medo de se tornar algo que não é muito bom -- porque o medo e o
desejo andam juntos.

No início, este caminho pode parecer um tanto desesperante. Às vezes, as
tendências e hábitos de uma vida inteira para se promover e evidenciar como uma
personalidade individual, são tão fortes, que sentimos que não devíamos ser
assim -- devíamos tentar ser ninguém.

Mas tentar ser ninguém ainda é ser alguém. O que estou a sugerir não é para se
tornar ninguém, mas para realizar a Verdade da mente. Então podemos permanecer
na Verdade, onde nos sentimos mais à vontade e em paz, em vez de neste ciclo
interminável de existência em que estamos sempre à procura de renascer de novo.
Em todos os níveis de existência, nunca encontraremos contentamento. Nunca nos
satisfazem, nem mesmo o melhor deles. Os estados condicionados de maior
felicidade, os \emph{jhānas}, ainda são insatisfatórios para nós. O Buddha
deixou claro que todas as formas de felicidade humana e sucesso mundano, são na
verdade terrivelmente decepcionantes, porque só nos podem gratificar
temporariamente. E no momento em que essa gratificação se vai, somos apanhados
no mesmo processo de buscar o renascer de novo, para nos tornarmos outra coisa,
para encontrar outro momento de felicidade. A vida torna-se tão cansativa.

Para viver com um corpo e atitude certa, comecemos a aceitá-lo como ele é e tudo
o que possa ter de correcto e errado em si, seja novo ou velho, homem ou mulher,
forte ou fraco. Este é o caminho para a verdadeira paz. Não nos procuremos
identificar com o nosso corpo ou tentemos transformá-lo noutra coisa. Quando
conhecemos a Verdade, então, nos momentos certos, podemos ser especiais de
acordo com o tempo e o lugar, sem que isso se torne um apego. Sentimos que
podemos manifestar-nos e desaparecer de acordo com o que for necessário. Não
estou a dizer que alguém tem que ficar entre as árvores pelo resto da sua vida.
Pode se ser algo útil e uma ajuda para os outros -- mas já não é um papel
permanente que se tenta manter e defender. Assim, começamos a sentir uma
sensação de liberdade e bem-estar.

Quando eu era jovem, era muito inibido -- dizer algo em público era
absolutamente assustador para mim. Mesmo quando eu estava na Marinha, só o facto
de ter que levantar a minha voz para dizer em público ``Sim, senhor!'' numa
lista de chamada, deixar-me-ia a tremer de inibição. Depois, tornei-me professor
primário. Ensinei crianças chinesas de 8 a 9 anos no Bornéu do Norte por dois
anos; não era assim uma grande ameaça. Mas depois, tornar-me monge na Tailândia
e, eventualmente ter que dar palestras para tailandeses em Tailandês\ldots{}! Toda
essa inibição veio ao de cima: as alturas que alcançava quando sentia que tinha
dado uma boa palestra e todas as pessoas diziam: ``Você é muito bom, Sumedho,
você pode dar um bom Dhamma.'' Então de vez em quando dava uma palestra de
Dhamma mesmo estúpida e pensava: ``Não quero dar outra palestra, nunca mais. Eu
não me tornei um monge para dar palestras.''

Mas a ideia era continuar a observar isso. Luang Por\footnote{\emph{Luang Por:} O termo
  tailandês que se traduz como ``Venerável Pai'', embora o Inglês não transmita a
  mistura de afecto e respeito que significa. É usado para se dirigir a um monge
  idoso.} Chah sempre me incentivou a ter consciência do orgulho, da vaidade, da
vergonha e da inibição que eu sentia. E, felizmente, na Tailândia, as pessoas
são de tal maneira que simplesmente ficam gratas por um monge dar uma palestra.
Mesmo não sendo uma palestra muito boa, parece não incomodá-los muito. Eles
ainda demonstram ficarem muito gratos por isso. Isso torna tudo muito mais
fácil. Uma vez, numa cerimónia da Kathina\footnote{\emph{Kathina:} A cerimônia
  realizada no final de um Vassa em que os leigos fazem ofertas ao mosteiro.} em
que tivemos que ficar acordados a noite inteira, Ajahn Chah disse: ``Sumedho,
vais dar uma palestra durante três horas esta noite.'' E eu só tinha falado meia
hora até àquele momento. Foi um stress, três horas!! E ele sabia.

Com Ajahn Chah, eu sentia que fosse o que fosse que ele dissesse, eu o faria
sempre. Então, sentei-me no cadeirão alto e falei por três horas. Tive que me
sentar ali e assistir às pessoas a levantarem-se e a irem-se embora; e tive que
me sentar ali e até assistir a pessoas a deitarem-se no chão e a dormirem na
minha frente. E no final das três horas, ainda lá havia algumas senhoras
educadas que permaneciam sentadas!

Aquilo não era Ajahn Chah a dizer, ``Ok Sumedho, vá lá e anime-os com algo
rejubilante! Divirta-os, sensibilize-os a sério!'' Comecei a perceber que o que
ele queria que eu fizesse, era ser capaz de olhar para este egocentrismo, a
pose, o orgulho, a presunção, o resmungo, a preguiça, o
``não-querer-me-chatear'', o querer agradar, o querer entreter, o querer obter
aprovação. Tudo isto surgiu ao longo das palestras dos últimos quinze anos. Mas
a meditação é aquilo em que cada vez mais uma pessoa sente uma compreensão
verdadeira do sofrimento da visão egocêntrica. E então, por meio desta
compreensão lúcida, realizamos a permanência no vazio.

Sempre que Ajahn Chah dava uma palestra, ele sentava-se, fechava os olhos e
então começava a falar -- e o que proferia seria apropriado para o momento e
lugar. Ele disse para nunca preparar uma palestra -- não se importava se eram
interessantes ou não -- para deixá-las simplesmente surgir. E quando existe
não-existência, não há mais eu, não há nenhum dos problemas que construímos com
``o que é que as pessoas pensam de mim? O que as pessoas dizem sobre mim?'' Ou a
rebelião -- ``Eles podem pensar o que quiserem, eu não me importo!'' (Mas não
pensa isso realmente, pensa? Caso contrário, não o diria certo?)

Às vezes, as personalidades manifestam-se nos momentos apropriados. À medida que
falamos, manifestamos a nossa personalidade. Ora, talvez ainda estejamos presos
à ideia de sermos uma pessoa na nossa mente. Mas essas são apenas condições que
surgem e cessam, e surgem por medo e desejo. Quando há vazio, a personalidade
ainda opera -- isso não significa que somos exactamente iguais, como as abelhas
numa colmeia. Ainda existe uma miríade de diferenças de carácter e personalidade
que se podem manifestar como carismáticas ou o que seja. Mas não há ilusão sobre
estas -- não há sofrimento.

Por exemplo, quando Ajahn Chah visitou pela primeira vez a Inglaterra, ele foi
convidado para uma refeição vegetariana na casa de certa senhora. Ela obviamente
esforçou-se imenso para confeccionar os mais deliciosos tipos de comida. Estava
ansiosa por oferecer esta comida e apresentava-se muito entusiasmada. Ajahn Chah
estava sentado a apreciar a situação e então de repente disse ``Esta é a
refeição mais deliciosa e maravilhosa que já alguma vez tive!''

Este comentário foi na realidade incrível, porque, na Tailândia, os monges não
devem fazer comentários sobre a comida. E no entanto, Luang Por de repente,
manifestou esta personalidade carismática que elogiava uma mulher que precisava
de ser elogiada, porque isso a fazia sentir muito feliz. Ele sentia a hora e o
lugar, a pessoa com quem estava, o que seria amável. Portanto, ele podia sair do
papel supostamente designado de acordo com a tradição e manifestar-se de
maneiras que fossem apropriadas.

Isso demonstra sabedoria e capacidade de responder a uma situação; não estar
rigidamente preso a uma convenção que o cega. Isto foi uma manifestação e
cessação ao mesmo tempo, porque nunca mais o ouvi dizer aquilo.

A mente vazia mantém-se tranquila, onde não há eu, nem algum medo ou desejo com
que se iludir. E ainda existe a capacidade de responder com compaixão e bondade
à situação presente de uma maneira adequada. É estranho, não é? Compare-se o
objectivo de Nibbāna, de não-existência, com o de tornar-se a melhor pessoa do
mundo, a mais forte ou a mais bela. Os valores mundanos envolvem possuir poder,
beleza, riqueza -- mas todos têm seus opostos, certo? O sucesso está sempre
ligado ao fracasso, a felicidade está sempre ligada à infelicidade, o elogio
está sempre ligado à censura. Boa sorte com má sorte. Portanto, se escolhermos
os valores mundanos de riqueza, poder, sucesso e elogio, levamos os outros junto
com eles porque são como os dois lados da mesma moeda. Não podemos separar um do
outro. Os valores mundanos nunca vão realmente permitir que nos sintamos à
vontade.

O mundo é um lugar inseguro. Não é pacífico. E não é onde realmente pertencemos.
Nós só começamos a compreender e a perceber a paz através do vazio, da
não-existência, do não-eu. E isso não é aniquilação, mas iluminação, liberdade,
verdadeira paz. Conhecimento verdadeiro.

\photoPage{image-009-anagarika-anjali.jpg}{Anagārika no cântico da noite}{Foto por Dick Todd}

\textbf{Reflexões sobre a Partilha de Bênçãos}

\begin{quote}
\emph{Cantemos agora as Reflexões sobre a Partilha de Bênçãos}

Através do bem que resulta da minha prática,

Que os meus mestres e guias espirituais de grande virtude,

A minha mãe, o meu pai e os meus familiares,

O Sol e a Lua, e todos os líderes virtuosos do mundo,

Que os Deuses mais elevados e as forças do mal,

Seres celestiais, espíritos guardiões da Terra e o Senhor da Morte,

Aqueles que são amigáveis, indiferentes ou hostis,

Que todos os seres recebam as bênçãos da minha vida.

Que brevemente cheguem à Tripla Bênção, e superem a morte.

Através do bem que resulta da minha prática,

E através desta partilha,

Que todos os desejos e apegos rapidamente cessem

Assim como os estados prejudiciais da mente.

Até realizar o Nibbana,

Em qualquer tipo de nascimento, que eu tenha uma mente justa,

Com consciência e sabedoria, austeridade e vigor.

Que as forças ilusórias não controlem,

nem enfraqueçam a minha decisão.

O Buddha é o meu excelente refúgio,

Insuperável é a proteção do Dhamma,

O Buddha solitário é o meu Nobre exemplo,

O Saṅgha é o meu maior suporte.

Que através desta supremacia

Desapareçam a escuridão e a ilusão.
\end{quote}

